% 2026-10-20用户反馈复核；图号、参考方向、数据与本题解答同步。
\documentclass[tikz,border=6pt]{standalone}
\usepackage{fontspec,xeCJK,amsmath,amssymb}
\IfFontExistsTF{Times New Roman}{\setmainfont{Times New Roman}}{\setmainfont{TeX Gyre Termes}}
\IfFontExistsTF{Heiti SC}{\setCJKmainfont[AutoFakeBold=1.5]{Heiti SC}}{\setCJKmainfont[AutoFakeBold=1.5]{FandolHei-Regular}}
\usetikzlibrary{arrows.meta,calc}
\definecolor{ink}{HTML}{243745}
\definecolor{blue}{HTML}{275DAB}
\definecolor{green}{HTML}{177C63}
\definecolor{red}{HTML}{BD4935}
\definecolor{orange}{HTML}{AA6C16}
\definecolor{purple}{HTML}{715393}
\begin{document}
\begin{tikzpicture}[x=1cm,y=1cm,font=\fontsize{9.3}{12.5}\selectfont,
v/.style={-{Latex[length=2mm,width=1.3mm]},line width=.85pt},
axis/.style={-{Latex[length=1.8mm,width=1.1mm]},ink!70,line width=.6pt},
aux/.style={ink!40,dashed,line width=.5pt},
box/.style={draw=ink!40,rounded corners=2pt,inner sep=5pt,align=center},
lab/.style={fill=white,inner sep=2pt}]
\path[use as bounding box] (-8,-3.6) rectangle(8,3.65);
\node[font=\bfseries\fontsize{10.5}{14}\selectfont,text width=15cm,align=center] at (0,3.23) {电感电流的连续条件与理想冲激};
\begin{scope}[shift={(-4,0)}]
\node[font=\bfseries,text=blue] at (0,2.55) {有界输入：状态连续};
\node[text width=7.2cm,align=center] at (0,1.5) {$L[i_L(t_0^+)-i_L(t_0^-)]=\int_{t_0^-}^{t_0^+}u\,dt$};
\draw[axis] (-2.55,-1.8)--(2.55,-1.8);
\draw[axis] (-2.45,-1.85)--(-2.45,0.55);
\node[anchor=west] at (-2.15,0.65) {$i_L$};
\node[] at (2.6,-2.0) {$t$};
\draw[blue,line width=1pt] (-2.2,-1.35)--(0,-0.9);
\draw[blue,line width=1pt] (0,-0.9)--(2.2,0.25);
\draw[aux] (0,-1.8)--(0,0.4);
\node[] at (0,-2.07) {$t_0$};
\node[] at (0,-2.6) {有界电压使积分趋零；斜率可改变};
\end{scope}
\begin{scope}[shift={(4,0)}]
\node[font=\bfseries,text=red] at (0,2.55) {若状态跃变：必须有冲激};
\node[] at (0,1.65) {$i_L(t_0^+)-i_L(t_0^-)=\Delta i\ne0$};
\draw[axis] (-2.55,-1.8)--(2.55,-1.8);
\draw[axis] (-2.45,-1.85)--(-2.45,0.55);
\node[anchor=west] at (-2.15,0.65) {$i_L$};
\node[] at (2.6,-2.0) {$t$};
\draw[red,line width=1pt] (-2.2,-1.3)--(0,-1.3);
\draw[red,line width=1pt] (0,0.1)--(2.2,0.1);
\draw[aux] (0,-1.3)--(0,0.1);
\node[] at (0,-2.07) {$t_0$};
\node[text=red] at (0,-2.6) {$u_{\rm imp}=L\Delta i\,\delta(t-t_0)$};
\end{scope}
\node[text width=15cm,align=center] at (0,-3.12) {连续性的条件是输入在缩小时间窗口内没有有限面积的冲激；理想源直接切换须另行核对模型。};
\end{tikzpicture}\end{document}